% ============================================================
%  2. Related Work
% ============================================================
\section{Related Work}
\label{sec:related}

\subsection{Transformer-Based NER}

BERT~\citep{devlin2019bert} established a new state of the art on CoNLL-2003
NER at 92.8~F1 by fine-tuning the full encoder with a linear classification
head.  Subsequent models---RoBERTa~\citep{liu2019roberta},
ALBERT~\citep{lan2020albert}, and SpanBERT~\citep{joshi2020spanbert}---improved
further.  All of these models share the fundamental limitation that their
subword tokenizers impose a fixed vocabulary and produce unpredictable
fragmentation of out-of-vocabulary (OOV) or misspelled tokens.

\subsection{Character-Level and Hybrid Models}

Early neural NER systems built word representations from character
CNNs~\citep{ma2016end} or BiLSTMs~\citep{lample2016neural}, achieving strong
OOV handling but lacking contextual pretraining.  CharBERT~\citep{ma2020charbert}
augments BERT's token embeddings with a parallel character-interaction channel
while sharing the same WordPiece tokenizer, improving robustness at the cost
of doubling the effective parameter count.

ByT5~\citep{xue2022byt5} and CANINE~\citep{clark2022canine} operate entirely
at the byte or character level, demonstrating that vocabulary-free encoding is
feasible in large Transformer models.  However, byte-level models require
significantly more compute per word due to the much longer input sequences.
\LightRet{} uses a \emph{pretrained} character embedder (\RetVec{}) rather
than learning character representations from scratch, achieving comparable
representational quality at far lower parameter count.

\subsection{Robustness to Noisy Text}

\citet{belinkov2018synthetic} showed systematic degradation of neural MT
models under character-level perturbations.  \citet{pruthi2019combating}
proposed a word-recognition model inserted before the NLP model as a defence,
but the subword tokenizer boundary remains a bottleneck.  Data-augmentation
approaches~\citep{wei2019eda,zhu2020freelb} fine-tune existing subword models
on perturbed text; \LightRet{} avoids this limitation entirely by operating at
the word level throughout.

\subsection{Knowledge Distillation for NLP}

DistilBERT~\citep{sanh2019distilbert}, TinyBERT~\citep{jiao2020tinybert}, and
MobileBERT~\citep{sun2020mobilebert} distil BERT into smaller Transformer
students, reducing parameters by 40--60\% with minimal accuracy loss.  All
remain subword-based.  Our distillation chain is distinctive: the final student
uses a completely different input representation than the teacher, and
distillation is performed in two sequential stages (sentence-level, then
token-level) to bridge the representational gap progressively.

\subsection{RetVec}

\RetVec{}~\citep{sander2023retvec} is a compact multilingual text vectorizer
trained via contrastive learning to produce similar embeddings for
typographically related word variants (misspellings, emoji representations,
Unicode lookalikes).  Originally developed for spam and content-moderation
classification at Google, \RetVec{} outputs a 256-dimensional word vector in
constant time regardless of word length.  To our knowledge, \LightRet{} is the
first system to use \RetVec{} as the backbone for a sequence labeling (NER)
task via structured knowledge distillation.
